%! Graphing Quadratic Functions \& Transformations — Unit 2, Lesson 2.1 — Lesson Plan
% GRADUAL-RELEASE plan (warm-up / guided notes & practice / debrief / close & start homework) on
% the Main Ideas / Notes table shape: the notes are one two-column table, four instruction rows
% then a Guided Practice row. THERE IS NO GROUP ACTIVITY — the guided notes carry the release row
% by row: I do (the definition, the display, the first problem of each grid) -> we do (the rest of
% each grid, then the Guided Practice row). THERE IS NO INDEPENDENT PRACTICE SET IN THE NOTES —
% the homework is the individual practice, and the period ends with students starting it. No
% tiers, no exit ticket, no extension box. Teacher-facing, 10pt; every student component is 12pt.
% Timing 5 / 35 / 10 / 10 (the SAAR allocation for the same 60-minute period).
\documentclass[10pt]{article}
\usepackage{algebra2-article}
\usepackage{algebra2-boxes}
\usepackage{graphicx}   % -article does NOT load graphicx; needed for thumbnails/screenshots
\graphicspath{{images/}}
\newcommand{\TallMath}[1]{$\displaystyle #1\rule[-1.4em]{0pt}{3.2em}$}
\newcommand{\CourseName}{Algebra 2}
\newcommand{\MeetingLength}{60 minutes}
\newcommand{\UnitNumberName}{Unit 2: Quadratic Functions \quad}
\newcommand{\LessonNumberName}{Lesson 2.1: Graphing Quadratic Functions \& Transformations}

\begin{document}
\begin{center}
    {\Large\bfseries \CourseName \\
    {\small\bfseries \UnitNumberName \LessonNumberName}}
\end{center}

\begin{tcolorbox}[colback=forestbg, colframe=forest, boxrule=0.9pt, arc=2mm,
    left=2mm, right=2mm, top=1.5mm, bottom=1.5mm]
    \textbf{Primary Objective:} Lesson~2.0 taught students to \emph{read} a parabola someone
    else drew. Here they learn to \emph{produce} one from its equation. They recognise the
    \textbf{three forms} of a quadratic --- \textbf{vertex form} $a(x-h)^2+k$, \textbf{standard
    form} $ax^2+bx+c$, and \textbf{intercept (factored) form} $a(x-p)(x-q)$ --- and read the
    \textbf{vertex}, the \textbf{axis of symmetry}, the \textbf{zeros} and the $y$-intercept
    straight out of whichever form is in front of them. They describe each as a
    \textbf{transformation} of the \textbf{parent function} $y=x^2$: a vertical shift from $k$, a
    horizontal shift from $h$, a stretch or squash from $|a|$, and a reflection when $a<0$ (so
    the vertex becomes a \emph{maximum}). \textbf{2.1 does no factoring and no solving} --- 2.2
    factors and 2.3--2.5 solve; here a factored form is \emph{given} and read.
    \\[2pt]
    \textbf{Standards (2023 VA SOL):} \textbf{A.F.2b, A.F.2c, A.F.2d} --- the Algebra~1
    quadratic-graphing cluster, reactivated as the foundation of this unit: investigate,
    analyze and describe the graph of a quadratic function, including its \emph{zeros /
    $x$-intercepts} and $y$-intercept (\textbf{A.F.2b}), its \emph{vertex} as a maximum or
    minimum (\textbf{A.F.2c}), and its \emph{axis of symmetry} (\textbf{A.F.2d}). Extended in
    Algebra~2 by \textbf{A2.F.1} --- for any function family, describe the transformations of
    the parent graph, including the \emph{horizontal shift} $f(x-h)$, the vertical shift, the
    stretch/compression, and the reflection --- applied here to the quadratic parent $y=x^2$;
    with characteristics investigated and analyzed per \textbf{A2.F.2a} (domain, range, zeros,
    intercepts) and \textbf{A2.F.2d} (extrema: the vertex as maximum or minimum).
    \\[2pt]
    \textbf{Lesson model:} \emph{Gradual release --- I do, we do, you do --- all of it inside the
    guided notes.} There is no group activity. The notes name the vocabulary as they teach it and
    deliver the instruction in \textbf{one Main Ideas / Notes table} on one context (the unit
    anchor from 2.0, $x^2-2x-3=(x-1)^2-4=(x+1)(x-3)$ --- one curve in three costumes): four
    instruction rows --- \emph{Its Shifts}, \emph{Stretch and Flip}, \emph{Vertex Form Reads
    Itself}, \emph{Three Costumes, One Curve} --- each a definition, a display, and a grid of
    problems (1--10) where you work the first and the class works the rest; then the Guided
    Practice row (11--14) is worked together with students holding the pen;
    \textbf{the homework is the individual practice}, and students start it in class while you
    circulate. The whole-class debrief is \emph{spoken} --- there is no practice set and no
    objective box on the notes page; the cover carries the targets.
\end{tcolorbox}

\renewcommand{\arraystretch}{1.4}
\boxguard
\begin{skillbox}[Priority Ideas \& Skills]{goldbox}
\begin{tabularx}{\linewidth}{@{}X|X@{}}
\textbf{Priority Ideas \& Skills} & \textbf{Key Understandings (the ``why'')} \\[2pt]
\hline
\vspace{-6pt}
\begin{itemize}[leftmargin=1.1em, nosep, after=\vspace{2pt}]
  \item Name the \textbf{parent function} $y=x^2$ and describe a parabola as a
        \textbf{transformation} of it.
  \item $f(x)+k$ shifts \emph{vertically}; $f(x-h)$ shifts \emph{horizontally} --- and
        subtracting inside moves it \textbf{right}.
  \item Read $a$: $|a|>1$ narrower, $0<|a|<1$ wider, $a<0$ opens down (vertex becomes a maximum).
  \item \textbf{Vertex form} $a(x-h)^2+k$: vertex $(h,k)$, axis $x=h$, with no work.
  \item \textbf{Standard form} $ax^2+bx+c$: $y$-intercept $(0,c)$; axis $x=-\frac{b}{2a}$; then
        substitute back for the vertex.
  \item \textbf{Intercept form} $a(x-p)(x-q)$: zeros $p$ and $q$; axis at the midpoint
        $x=\frac{p+q}{2}$.
  \item Show that all three forms of one curve give the \emph{same} vertex and axis.
\end{itemize}
&
\vspace{-6pt}
\begin{itemize}[leftmargin=1.1em, nosep, after=\vspace{2pt}]
  \item A quadratic's \emph{form} is a costume, not a different curve. Choosing a form is
        choosing which feature you want handed to you for free.
  \item \textbf{Target misconception: the horizontal-shift sign flip.} Students read $(x+4)^2$
        as ``right~4.'' It is \emph{left}~4: $y=(x-1)^2$ moves the parent \textbf{right} $1$
        while $y=x^2-1$ moves it \textbf{down} $1$ --- inside versus outside. Sharpen it against
        intercept form, where the signs flip the \emph{other} direction: in $(x+1)(x-3)$ the
        zeros are $-1$ and $3$. \textbf{The same bracket does opposite jobs in the two forms, and
        which form you are in decides.} Notes problems~9 and~10 settle it; Guided Practice
        problem~14 tests it on new ground.
  \item ``Narrower'' is a claim about \emph{outputs}, not about how the picture looks --- notes
        problem~3 checks it in a table ($x^2$ gives $4$ at $x=2$; $3x^2$ gives $12$).
  \item The axis of symmetry is the same line however you compute it, because symmetry is a
        property of the curve, not of the algebra.
\end{itemize}
\end{tabularx}
\end{skillbox}

\boxguard
\begin{skillbox}[{Vocabulary, Concepts, \& Theorems --- taught directly in the guided notes}]{sky}
\renewcommand{\arraystretch}{1.3}
\begin{tabularx}{\linewidth}{@{}l X@{}}
\textbf{Parent function $y=x^2$} & The simplest quadratic; every other parabola is a transformation of it. \\
\textbf{Transformation} & A change to the parent's rule that shifts, stretches or squashes, or reflects its graph. \\
\textbf{Vertex form $a(x-h)^2+k$} & Hands you the vertex $(h,k)$ and the axis $x=h$ with no work --- flip the sign inside the bracket. \\
\textbf{Standard form $ax^2+bx+c$} & Hands you the $y$-intercept $(0,c)$; the axis is $x=-\frac{b}{2a}$, and that $x$ substituted back gives the vertex. \\
\textbf{Intercept (factored) form $a(x-p)(x-q)$} & Hands you the zeros $p$ and $q$; the axis sits at their midpoint $x=\frac{p+q}{2}$. \\
\textbf{Axis of symmetry} & The vertical line through the vertex; the parabola mirrors across it, so points pair off at equal heights. \\
\end{tabularx}
\end{skillbox}

\begin{fixedskillbox}[Lesson at a Glance --- \MeetingLength]{forestbg}
\renewcommand{\arraystretch}{1.25}
\begin{tabularx}{\linewidth}{@{}l c X X@{}}
\textbf{Phase} & \textbf{Min} & \textbf{Students} & \textbf{Teacher} \\
\hline
Warm-Up & 5 & Three prior skills, alone, then check with a neighbour & Circulate; debrief item~3
  aloud (``$y=(x-3)^2$ moved \emph{right} $3$ --- with a minus sign'') and leave it hanging \\
Guided Notes \& Practice & 35 & Fill the vocabulary box and the notes table row by row as each
  term is named, working the rest of each grid themselves; then Guided Practice (11--14)
  together, pen in their hand & \textbf{I do} the hook, then each row's definition, display, and
  first problem (1, 3, 6, 9) (20) $\to$ \textbf{we do} the rest of each grid, then the Guided
  Practice row (15): one new curve, read three ways \\
Debrief & 10 & Pens down, papers up --- answer aloud, nothing to fill in & Put the three costumes
  back on the board; take the ``$(x+4)^2$ shifts right'' reflex and the ``zeros are $1$ and
  $-5$'' error \\
Close \& Start Homework & 10 & \textbf{Start the homework in class}, alone and silent & Launch
  item~1(a) aloud, then circulate for your formative read; preview Lesson~2.2; the rest is due
  the first class after two study halls \\
\end{tabularx}
\end{fixedskillbox}

\begin{fixedskillbox}[Warm-Up --- Activate Prior Knowledge (5 min)]{forestbg}
\begin{minipage}[t]{0.48\linewidth}
    \textbf{The three items and what each seeds}
    \begin{itemize}
        \item \textbf{1.} Expand $(x-1)^2-4$ and $(x+1)(x-3)$; both give $x^2-2x-3$.
              \textbf{Seeds:} FOIL, the one mechanical skill today needs --- and it quietly
              plants the hook: the two expressions are the \emph{same} curve.
        \item \textbf{2.} Transformations of $y=|x|$ from Lesson~1.3 --- up/down, left/right, the
              sign flip. \textbf{Seeds:} the whole transformation vocabulary, already in the
              room before the parabola gets it.
        \item \textbf{3.} The parent $y=x^2$ beside $y=(x-3)^2$, pre-drawn: it moved
              \emph{right}~$3$. \emph{Does that surprise you, given the minus sign?}
              \textbf{Seeds:} the target misconception, which notes row~1 answers.
    \end{itemize}
\end{minipage}
\hfill
\begin{minipage}[t]{0.48\linewidth}
    \textbf{Running it}
    \begin{itemize}
        \item \textbf{Debrief item~3 aloud, briefly}, and land only this: \emph{it moved right,
              and the sign inside says minus --- hold on to that.} Do not explain it yet; the
              hook's vote is worth more than the thirty seconds.
        \item Item~1's punchline --- ``two different-looking expressions, one answer'' --- is
              worth five extra seconds, because the hook opens on exactly that picture.
        \item If it runs long, cut item~2's third bullet and keep item~3 whole.
        \item The page is 12pt like the rest of the packet and fits one side. If you add to it,
              check that it still does.
    \end{itemize}
\end{minipage}
\end{fixedskillbox}

\boxguard
\begin{skillbox}[Guided Notes \& Practice --- I do, then we do (35 min)]{forestbg}
\textbf{This is the whole lesson.} There is no group activity, and there is no independent
practice set on this page: \textbf{the homework is the individual practice}, started in class.
\begin{multicols}{2}
\textbf{Hook (3 min, page 1).} Put the three expressions up, stacked, beside \emph{one}
pre-drawn parabola: $x^2-2x-3=(x-1)^2-4=(x+1)(x-3)$. Then the vote that carries the lesson:
\textbf{``Three different parabolas, or one parabola wearing three costumes?''} Take a hand
count, hear one reason for each, and \emph{do not settle it} --- row~4 does. Then the two
questions that shape the whole table: \emph{which costume tells you the vertex fastest? which
one the zeros?} Fill the vocabulary rows \emph{as you go}, never in advance.

\textbf{I do, row by row (about 17 min).} The notes are one table; each row is one idea. For
each row: read the definition aloud, mark up the display, work the \emph{first} problem of the
grid on the board, and hand the rest to the room.

\emph{Its Shifts (pages 1--2, 5 min).} The parent $y=x^2$; outside moves it up/down, inside
moves it left/right. Mark up the one grid carrying $y=x^2$, $y=x^2+3$ and $y=(x-2)^2$, and fill
the \texttt{labelbox} (``right~2''). Work \textbf{1} ($+3$ slides it up, vertex $(0,3)$); the
class works \textbf{2} --- the sign flip, stated for the first time.

\emph{Stretch and Flip (page 2, 4 min).} $a$ does two jobs: sign is direction, size is width.
\textbf{Work 3 --- the value table --- and do not skip it:} ``narrower'' is a claim about
outputs ($4$ against $12$ at $x=2$), not about what the picture looks like. The class works
\textbf{4} and \textbf{5}.

\emph{Vertex Form Reads Itself (page 3, 4 min).} The five \texttt{stepnum} steps, then the
anchor $(x-1)^2-4$: $a=1$, $h=1$ (\emph{flip the sign inside}), $k=-4$. Work \textbf{6}; the
class works \textbf{7} and \textbf{8} (8 runs the skill backwards: vertex to rule).

\emph{Three Costumes, One Curve (pages 3--4, 5 min) is the lesson.} Standard form gives $(0,-3)$
free; the \texttt{work} block gets the axis $x=-\frac{-2}{2(1)}=1$ and substitutes back to
$y=-4$; intercept form gives zeros $-1$ and $3$ and the midpoint $\frac{-1+3}{2}=1$.
\textbf{Say it plainly: all three agree, because it is one curve} --- and the hook's vote is
settled. \textbf{9 is the crux:} let someone defend ``$(x+4)^2$ shifts right~4'' and let the room
catch it. \textbf{10 is the crux sharpened} --- the same bracket in the two forms, running
\emph{opposite} ways.

\columnbreak
\textbf{We do (about 15 min, page 4).} The Guided Practice row: $y=-(x-2)^2+9=-x^2+4x+5=
-(x+1)(x-5)$, pre-drawn. Students hold the pen; you hold the questions. \textbf{11:} vertex
$(2,9)$, axis $x=2$, opens down, so $9$ is a \emph{maximum}. \textbf{12:} $c=5$ gives $(0,5)$;
the \texttt{work} block gives $x=-\frac{4}{2(-1)}=2$ and $y=9$. \textbf{13:} zeros $-1$ and $5$,
midpoint $2$. \textbf{14 is the crux transferred and the problem to protect:} all three axes are
$x=2$, and they \emph{must} agree because it is one curve; the classmate who read the zeros as
$1$ and $-5$ flipped the signs backwards --- $(x+1)$ is zero at $x=-1$.

\textbf{The pen must actually change hands.} Guided Practice is the only place today where you
can watch a student decide, so do not narrate it. Put problem~13's zeros on them cold, and take
a hand count on ``do the three axes agree?'' before anyone speaks.

Circulate during Guided Practice and demand the reason, not the word:
\begin{itemize}[leftmargin=1.1em, nosep]
  \item \textbf{Problem 11:} ``What is $h$ here --- and did you flip the sign inside?''
  \item \textbf{Problem 13:} ``$(x+1)$ is zero at which $x$? Put it in and check.''
  \item \textbf{Problem 14, the crux transferred:} ``You got $x=2$ three times. Why did that
        \emph{have} to happen?''
\end{itemize}
While circulating, pick two papers to put up in the debrief --- one problem~14 with all three
axes agreeing and the zeros named correctly, and one that read the zeros as $1$ and $-5$.
\textbf{Your formative read comes twice}: here, and again in the last ten minutes as students
start the homework.
\end{multicols}
\end{skillbox}

\boxguard
\begin{skillbox}[Debrief --- whole class, spoken (10 min)]{forestbg}
Pens down, papers up. \textbf{This debrief is spoken} --- there is no box at the end of the
notes to fill in, so it lives entirely on the board and in the room.
\begin{multicols}{2}
\begin{enumerate}[leftmargin=1.3em, nosep, itemsep=3pt]
  \item \textbf{Put the three costumes back up} --- $x^2-2x-3$, $(x-1)^2-4$, $(x+1)(x-3)$ --- and
        make a student say the vote's answer: \emph{one curve}. Then have the room name which
        form hands you which feature, without looking.
  \item \textbf{The sign flip.} $y=(x-1)^2$ moves the parent \emph{right}~$1$; $y=x^2-1$ moves it
        \emph{down}~$1$. Inside versus outside. Take the ``$(x+4)^2$ shifts right~4'' reflex here
        and kill it aloud: $x+4$ is $x-(-4)$, so $h=-4$.
  \item Put up the two problem~14s you picked. The one that read the zeros as $1$ and $-5$
        flipped the signs backwards --- make the room substitute $x=-1$ into $(x+1)(x-5)$ and
        watch it give zero. That is the moment ``the form decides what the bracket means'' stops
        being a slogan.
  \item Take ``$a$ only makes it skinny'' last: $a<0$ also flips it, and a flipped parabola's
        vertex is a \emph{maximum} --- Lesson~2.0's word, now read off the equation.
\end{enumerate}
\columnbreak
\textbf{Two things to demand aloud.} First, \textbf{make a student read a vertex form without the
notes} --- ``$y=(x-1)^2-4$: vertex $(1,-4)$, axis $x=1$, opens up, and I flipped the sign inside
to get $h=1$.'' Second, \textbf{make them say why three different computations gave the same
axis} --- because symmetry belongs to the curve, not to the algebra.

\textbf{If the debrief runs short}, cut item~4 and item~1 --- items~2 and~3 are the two that
cannot be cut. \textbf{Do not let the debrief eat the last ten minutes} --- the homework start is
not optional padding, it is the individual practice.
\end{multicols}
\end{skillbox}

\boxguard
\begin{skillbox}[Homework --- scored, started in class, due the first class after two study halls]{goldbox}
Six items spanning A.F.2b/c/d and A2.F.1 / A2.F.2a / A2.F.2d, two pages, no extension:
\textbf{1} three vertex-form rules read for vertex, axis, direction and width; \textbf{2} the
deliberate \emph{inside-versus-outside} contrast pair, $y=(x+3)^2$ against $y=x^2+3$, settled by
a table of values at $x=-3,0,3$ and closing with the ``why does subtracting inside move it the
other way?'' sentence; \textbf{3} \emph{reading a pre-drawn graph and then writing its rule in
vertex form} --- the only item that runs backwards, and the best single predictor of whether the
sign flip stuck; \textbf{4} intercept form $y=2(x-1)(x+5)$ --- zeros, midpoint axis, vertex by
substitution in a \texttt{work} block --- plus the special case $y=a(x-3)^2$, where the two zeros
coincide and there is only \emph{one} $x$-intercept; \textbf{5} the stone arch
$h(x)=-\frac14(x-8)^2+16$ with a \texttt{work} block for $h(4)$ and an interpret-the-answer
follow-up; \textbf{6} an SOL-style multiple-choice item on the direction of a horizontal shift,
the formative check.

\textbf{Start it in the last ten minutes of the period, in class and alone.} Item~1(a) is the
launch; item~2's table is where the misconception shows first, and item~3's rule is where it
shows fatally. \textbf{Item~6 carries the check when you score it:} (A) is the sign flip read
backwards (the target misconception), (C) and (D) treat an inside number as an outside one.

\textbf{Two pages is a ceiling at 12pt}, so use any extra items as debrief cold-calls rather
than a third page. \textbf{DeltaMath override:} transformations of quadratics and reading vertex
form \emph{are} well covered on DeltaMath --- ``Graphing Quadratic Functions in Vertex Form''
together with ``Transformations of Quadratic Functions'' is the swap --- and that replaces this
page and strikes its score cell on the cover. Never assign both; DeltaMath is thinner on the
three-forms-agree idea, so keep items~4 and~6 if you do. \textbf{Preview:} Lesson~2.2,
\emph{Factoring Quadratic Expressions} --- where the factored form we read today actually comes
from.
\end{skillbox}

\boxguard
\begin{skillbox}[Watch For (while circulating)]{redbox}
\begin{itemize}[leftmargin=1.2em, nosep]
  \item \textbf{Reading $(x+4)^2$ as ``right~4''} (notes~2, 9 and~10; HW~1(b), 2, 6) --- the
        lesson's \textbf{target misconception}. Probe: \emph{``What $x$ makes the bracket zero?
        That is where the vertex went.''}
  \item \textbf{Flipping the signs the wrong way in intercept form} --- reading $(x+1)(x-5)$ as
        zeros $1$ and $-5$ (notes~10 and~14, the crux transferred; HW~4). Probe: \emph{``Put
        $x=-1$ in. Did you get zero?''}
  \item Confusing inside with outside --- treating $x^2-1$ as a horizontal shift (notes~1--2,
        HW~2). Probe: \emph{``Is the number inside the bracket or outside it?''}
  \item Saying ``narrower'' from the picture without checking outputs (notes~3, HW~1). Probe:
        \emph{``At $x=2$, what does each one give?''}
  \item Forgetting that $a<0$ makes the vertex a \emph{maximum}, not a minimum (notes~4 and~11,
        HW~1(b)--(c), 5).
  \item In HW~4(d), thinking a repeated zero means \emph{no} $x$-intercepts rather than exactly
        one, with the vertex sitting on the $x$-axis.
  \item Trying to \emph{factor} or \emph{solve} in notes~10 or HW~4 --- that is Lessons 2.2--2.3.
        Today the factored form is handed to you.
\end{itemize}
Cold-call prompts: ``Which form is this, and what does it hand you for free?'' ``What is $h$ ---
did you flip the sign?'' ``Inside or outside the bracket?'' ``$(x+p)$ is zero at which $x$?''
``Do your three axes agree --- and why must they?''
\end{skillbox}

\boxguard
\begin{skillbox}[Close \& Start Homework (10 min)]{goldbox}
\textbf{Students start the homework here, in class, alone and silent --- this is the individual
practice, and the first minutes of it are your best formative read.} Hand the launch first ---
six items on paper, scored, due the first class after two study halls (announce the date in class
and on TurtleNet) --- and work item~1(a) aloud in thirty seconds so nobody stalls on the page.
Then circulate on item~2's table and item~3's rule, sorting into three piles: flipped the sign
inside correctly and wrote $y=(x+2)^2+1$ (got it); wrote $y=(x-2)^2+1$ or called $(x+3)^2$ a
shift right (the target misconception, still alive); confused inside with outside entirely (send
them back to the warm-up's item~3 picture). Close on what changed: \emph{``You walked in able to
read a parabola someone drew for you. You walk out able to take the equation, name which of the
three costumes it is wearing, and say where its vertex, axis, zeros and $y$-intercept are ---
without drawing anything.''}
\end{skillbox}

% ── Teacher notes — one per component, in packet order (three: no activity, no practice set) ──
\begin{teachernote}[Warm-Up]
Item~1 is the mechanical core and the quiet plant: both expansions land on $x^2-2x-3$, which is
exactly the hook's picture, so let the room notice it out loud before you move on. Item~2 is the
bridge from Lesson~1.3 --- the transformation words (shift, left/right, flip) are already theirs,
and today only the parent changes. Item~3 is the make-or-break setup for the whole lesson:
$y=(x-3)^2$ sits three units \emph{right} of the parent, with a minus sign inside. Debrief it in
one sentence --- \emph{``it moved right; the sign says minus; hold on to that''} --- and leave it
hanging. Do \emph{not} explain the mechanism here; notes row~1 pays it off within eight minutes,
and the hook's vote is a better use of the thirty seconds. The page is 12pt like every other
student component and fits one side, blank and keyed; both \texttt{work} blocks in item~1 are
byte-identical in the key.
\end{teachernote}

\begin{teachernote}[Guided Notes \& Practice]
\textbf{This one document is the lesson --- pace it 3 / 17 / 15, then 10 for the debrief, then 10
for the homework start.} Page~1 is the six vocabulary rows and the hook; fill the rows \emph{as
each term is named}, never in advance, and take the three-costumes vote without resolving it.
Then the table, row by row: read the definition, mark up the display, work the first problem of
the grid (\textbf{1, 3, 6, 9}), and hand the rest to the room.

\emph{Its Shifts} must move: five minutes on the parent, the one grid with three labelled curves,
the \texttt{labelbox}, and problems 1--2. \emph{Stretch and Flip} is four, and
\textbf{problem~3 --- the value table --- is the one to protect in that row}: it is the only place
all day where ``narrower'' is checked with numbers instead of asserted, and students who skip it
carry a picture-based idea of $a$ into Unit~3. \emph{Vertex Form Reads Itself} is four: the five
steps, then the anchor, then 6--8 (8 runs the skill backwards).

\textbf{Spend the time in \emph{Three Costumes, One Curve}, which carries the misconception.}
Standard form first ($c$ free, then the \texttt{work} block for the axis and the vertex), then
intercept form (zeros, then the midpoint), then the sentence that settles the hook: \emph{all
three agree, because it is one curve}. \textbf{9 is the crux:} let someone defend ``$(x+4)^2$
shifts right~4'' and let the room catch it --- do not resolve it by repeating the rule, ask what
$x$ makes the bracket zero. \textbf{10 is the crux sharpened, and the sentence to land:} the same
bracket shifts \emph{left} in vertex form and gives the zero $-4$ as a factor. Say ``which form
you are in decides'' and write it on the board.

\textbf{Guided Practice (11--14) is where the pen changes hands.} 11 is vertex form, 12 standard
(with the second \texttt{work} block), 13 intercept, and \textbf{14 the crux transferred} --- all
three axes are $x=2$, and the classmate who read the zeros as $1$ and $-5$ flipped the signs
backwards. If time is short, cut 12, never 14.

\textbf{There is no independent practice set on this page, and that is deliberate --- the
homework is the individual practice.} Guided Practice is the last row of the table; the period
ends with students starting the homework in front of you, which is where your second formative
read comes from. \textbf{2.1 does no factoring and no solving:} if a student starts factoring
$x^2-2x-3$, tell them that is Lesson~2.2 and hand the factored form back to them. Two
\texttt{work} blocks in the notes, byte-identical in blank and key; every other answer space is a
fixed-height grid cell.

\textbf{Debrief (10 min), spoken.} Three costumes back on the board, then the sign flip, then the
two problem~14s; items~2 and~3 are the two that cannot be cut. \textbf{Do not let the debrief eat
the last ten minutes.}
\end{teachernote}

\begin{teachernote}[Homework]
\textbf{Scored; due the first class after two study halls --- announce the date in class and on
TurtleNet.} The ten minutes at the end of the period are a \emph{start}, not a completion: put
students on 1 and 2 while you can still catch a wrong start, and let 3--6 go home. Item~1(c)'s
``same width'' is the cell most often left blank --- $a=-1$ changes direction only. Item~2 is the
contrast pair: expect the table to be right and part~(a) to be wrong, which is exactly the
diagnosis you want. \textbf{Item~3 is the one that runs backwards} --- vertex $(-2,1)$, axis
$x=-2$, $y$-intercept $(0,5)$, rule $y=(x+2)^2+1$; a student who writes $y=(x-2)^2+1$ has the
misconception still alive, and that is the single most useful thing you will see all day. Item~4
is intercept form read, not solved: zeros $1$ and $-5$, axis $x=-2$, vertex $(-2,-18)$ from the
\texttt{work} block, and part~(d)'s repeated zero gives exactly \emph{one} $x$-intercept with the
vertex on the $x$-axis. Item~5's arch: peak $16$ ft at $x=8$ straight out of vertex form, $16$ ft
wide at ground level, and $h(4)=12$ means the arch is $12$ ft high four feet from the left end.
\textbf{When you score the page, sort item~6 into three piles}: chose B (got it); chose A (the
sign flip read backwards --- the target misconception, still alive); chose C or D (inside treated
as outside). If more than a third land outside B, open Lesson~2.2 with a one-minute re-look at
``what $x$ makes the bracket zero'' rather than letting it ride. Two \texttt{work} blocks
(items~4c and~5c), byte-identical in the key. \textbf{DeltaMath override:} this content \emph{is}
well covered there --- ``Graphing Quadratic Functions in Vertex Form'' plus ``Transformations of
Quadratic Functions'' --- so a swap is legitimate here; it replaces this page and strikes its
score cell on the cover, and you should keep items~4 and~6 on paper if you swap, because the
three-forms-agree idea is thin on DeltaMath. Never assign both.
\end{teachernote}

\end{document}
